%% ===================================================================
\section{Conclusion}
\label{sec:conclusion}

We built reproducible ray-traced corpora of \nCities{} cities in two bands
for cross-city transfer of path-loss models and evaluated transfer where a
network planner uses it. Tile labels of the log-distance law are
ill-conditioned and, once their correlated errors are accounted for,
imprecise: the cluster-robust standard error of the exponent is \seRatio{}
times the naive one, and the exponent depends on the serving sites, the tile
size and the cut-off of the data. Transferring tile parameters with learned
regressors improves on the source median in \nGbdtMedian{} of \nFolds{}
cities, although not significantly after correction for multiple
comparisons; kernel transfer improves on it less, and our analysis of its limiting regimes and the gain
from learning its metric locate its weakness in the isotropic distance. At
the network level, even perfect tile parameters leave a path-loss error of
\SI{\netOracle}{\decibel}, assign only \assocOracle\% of the pixels to the
correct site and misplace cell edges by \SI{\edgeOracle}{\metre}; composing
the parameters of the tiles along each path improves on this with exact
labels, but not in path loss or serving site once they are transferred. Site-aware
hurdle link transfer, which predicts each link and whether it exists from its
obstruction geometry, reduces the path-loss error to \SI{\netLink}{\decibel},
raises serving-site accuracy to \assocLink\% and lowers the planning regret in
every held-out city, and a city-level certificate turns its accuracy into
coverage that can be promised at a stated risk.

Future work should validate these findings on measured data with known site
positions, extend the corpora to terrain, antenna patterns and further
bands, and study certificates that adapt to the target city from a small
survey.
